% LLM-MetaOpt -- paper skeleton.
% Writing rules: American English, passive voice, no semicolons, no em dashes,
% no bold or itemized lists inside main sections, external BibTeX file.
% Compile with: latexmk -pdf main.tex   (or upload the folder to Overleaf)

\documentclass[11pt,a4paper]{article}

\usepackage[utf8]{inputenc}
\usepackage[T1]{fontenc}
\usepackage{amsmath,amssymb}
\usepackage{graphicx}
\usepackage{booktabs}
\usepackage{url}
\usepackage[hidelinks]{hyperref}
\usepackage{algorithm}
\usepackage{algpseudocode}

\title{Language Models as Heuristic Meta-Optimizers and Intelligent Initializers for Variational Quantum Algorithms with Integrated Explainability}
\author{Daniel Mart\'in P\'erez\thanks{Data Science and Big Data Research Lab, Universidad Pablo de Olavide, Seville, Spain.}}
\date{\today}

\begin{document}
\maketitle

\begin{abstract}
Variational quantum algorithms are limited by non-convex optimization on noisy intermediate-scale hardware. Local optimizers stall on barren plateaus and on rough landscapes produced by physical noise, while calling a language model at every iteration is infeasible because of the input-output latency tax of hosted inference. This work studies a hybrid control architecture in which a language model supervises a fast local optimizer. The language model proposes the initial angles from the symbolic description of the problem and, every $N_w$ iterations, reads a telemetry window of the training run, reports a diagnostic label, justifies it and selects one macroscopic intervention. The intervention is limited to a scaling of the learning rate, a Gaussian perturbation of the variational angles and a restart. The study quantifies the effect of the controller on the final energy gap and separates it from the effect of perturbation alone with random, rule-based and learned classical controllers. Diagnosis accuracy is measured against simulator ground truth, and the latency budget of the slow loop is reported as a distribution. Experiments are carried out on Heisenberg, transverse-field Ising and XY models with two to four qubits under depolarizing noise of zero, two and five percent, using Qiskit and the statevector and density-matrix simulators.
\end{abstract}

\section{Introduction}
% Motivation: NISQ limitations, gradient decay, optimization bottlenecks.
% Gap: language models have been used to propose circuits and to act as
% black-box optimizers, yet their use as supervisory meta-controllers of a
% variational optimizer with an explicit explanation protocol is unexplored.
% Contributions are listed in prose (no itemize), for example: (i) the fast-slow
% loop formulation, (ii) the explainability protocol and its evaluation against
% ground truth, (iii) an open benchmark with counterfactual labels, (iv) the
% latency accounting.

\section{Related Work}
% Slots to fill with verified references (paper/refs.bib):
% 1. Variational quantum algorithms and barren plateaus.
% 2. Classical optimizers for VQAs and the parameter-shift rule.
% 3. Language models as optimizers and black-box search controllers.
% 4. Language models for quantum computing tasks.
% 5. Reinforcement learning and meta-learning for quantum control.
% 6. Explainability in quantum machine learning.

\section{Background}
\subsection{Variational quantum algorithms}
% Hamiltonian H, ansatz U(theta), expectation value E(theta), parameter-shift rule.
\begin{equation}
E(\vec{\theta}) = \langle \psi(\vec{\theta}) \rvert H \lvert \psi(\vec{\theta}) \rangle = \mathrm{Tr}\left(\rho(\vec{\theta}) H\right)
\end{equation}
\subsection{Barren plateaus and local minima}
\begin{equation}
\mathrm{Var}\left[\partial_{\theta_i} E(\vec{\theta})\right] \sim \mathcal{O}(2^{-n})
\end{equation}
\subsection{Language models as optimizers}
% OPRO-style prompting, evolutionary search with language models, and the
% latency argument for supervising rather than iterating.

\section{Method}
\subsection{Fast-slow loop}
% Algorithm 1: SPSA for N_w steps, query the controller, apply the action.
\begin{algorithm}
\caption{SPSA with a slow-loop controller}\label{alg:fastslow}
\begin{algorithmic}[1]
\Require initial angles $\vec{\theta}_0$, budget $T$, window $N_w$, controller $\mathcal{C}$
\State $\eta \gets 1$
\For{$k = 0$ \textbf{to} $T-1$}
  \If{$k \bmod N_w = 0$ \textbf{and} $k > 0$}
    \State $T_k \gets$ telemetry window of the last $N_w$ steps
    \State $a \gets \mathcal{C}(T_k)$
    \State $\vec{\theta}_k \gets \mathrm{intervene}(\vec{\theta}_k, a)$ and update $\eta$
  \EndIf
  \State $\vec{\theta}_{k+1} \gets \mathrm{SPSA\text{-}step}(\vec{\theta}_k, \eta)$
\EndFor
\end{algorithmic}
\end{algorithm}
\subsection{Telemetry window}
\begin{equation}
T_k = \left[E_{k-N_w \ldots k}, \lVert \nabla E \rVert_2, \eta_k, \mathrm{Var}(\vec{\theta})\right]
\end{equation}
\subsection{Explainability protocol}
% Diagnosis, justification, intervention. Mention the JSON schema constraint and
% that the reported label is compared with simulator ground truth.
\subsection{Intervention space}
\begin{equation}
\vec{\theta}_{k+1} = \vec{\theta}_k - \eta_{\mathrm{LLM}} \, \hat{g}_k + \delta \, \mathcal{N}(0, I)
\end{equation}

\section{Experimental setup}
% Problems, noise levels, seeds, conditions, statistics. Table 1 is generated by
% code.aggregate and can be included directly:
% \input{../results/tables/summary.tex}
\subsection{Benchmarks}
\subsection{Controllers and controls}
\subsection{Statistical analysis}

\section{Results}
\subsection{Optimization quality}
\subsection{Attribution to the controller}
\subsection{Diagnosis accuracy against ground truth}
\subsection{Latency and cost accounting}

\section{Discussion}
% When does the slow loop help, and when does it not. Sensitivity to N_w and to
% the model. Cost-effectiveness of a trained specialist against a general model.

\section{Limitations}
% Simulator-only evidence, small qubit counts, single-family prompts, model
% drift, prompt sensitivity, and the fact that the oracle controller is not
% deployable.

\section{Conclusion}

\bibliographystyle{plain}
\bibliography{refs}

\end{document}
